% Generated draft: requires mathematical and source review.
\documentclass[11pt]{article}
\usepackage[margin=25mm]{geometry}
\usepackage{fontspec}
\setmainfont{Latin Modern Roman}
\usepackage{amsmath,amssymb,mathtools,unicode-math}
\setmathfont{Latin Modern Math}
\usepackage{enumitem,needspace,xurl,hyperref}
\setlength{\parindent}{0pt}
\setlength{\parskip}{5pt}
\setlength{\emergencystretch}{3em}
\newcommand{\N}{\mathbb{N}}
\newcommand{\Z}{\mathbb{Z}}
\newcommand{\Q}{\mathbb{Q}}
\newcommand{\R}{\mathbb{R}}
\newcommand{\C}{\mathbb{C}}
\newcommand{\F}{\mathbb{F}}
\newcommand{\A}{\mathbb{A}}
\newcommand{\PP}{\mathbb{P}}
\newcommand{\E}{\mathbb{E}}
\newcommand{\eps}{\varepsilon}
\newcommand{\dd}{\,\mathrm{d}}
\begin{document}
\section*{UnsolvedMath ID: 9700038}
\textbf{Source problem code:} catalogue entry
\subsection{Definitions and mathematical statement}
Use the standard space of isomorphism classes of connected countable locally finite rooted graphs $(G,o)$, with the local topology determined by finite-radius rooted balls. A law $\mu$ is unimodular if it satisfies the mass-transport principle: for every nonnegative Borel function $F$ of a graph with an ordered pair of distinguished vertices,
\[
 \E_\mu\sum_{v\in V(G)}F(G,o,v)
     =\E_\mu\sum_{v\in V(G)}F(G,v,o).
\]
The equality is allowed in $[0,\infty]$.

For a finite graph $G_n$, choose $o_n$ uniformly among its vertices and view the connected component of $o_n$ as the rooted graph. A local weak limit $\mu$ means that, for every fixed radius $R$ and finite rooted radius-$R$ ball type $H$,
\[
 \Pr\bigl(B_R(G_n,o_n)\cong H\bigr)
       \longrightarrow
 \mu\{(G,o):B_R(G,o)\cong H\}.
\]
A law admitting such a finite-graph approximation is called sofic in this setting.

The Aldous--Lyons question is whether every unimodular law on countable locally finite rooted graphs is sofic. No uniform degree bound is imposed in the supplied statement. Connected rooted graphs and uniform finite roots are the standard interpretation made explicit here; marks, multiple edges, and other category variants are not specified.

The supplied research summary reports a negative resolution: a unimodular random graph exists that is not a local weak limit of any sequence of finite graphs. It attributes this to Bowen--Chapman--Lubotzky--Vidick (2024). This is a source-reported solution, not a construction or proof supplied by this draft.
\subsection{Short English statement}
Can every random rooted graph satisfying mass transport be approximated locally by finite graphs rooted at a uniform vertex? The supplied summary answers no, citing a 2024 nonsofic unimodular example. It also reports a 2024/2025 undecidability strengthening, whose input/encoding formulation is not included. The solved label is retained with that attribution, while the source provenance still requires a current primary-literature audit.
\subsection{Sources}
\begin{itemize}
\item Local UnsolvedMath registry, permanent ID 9700038, source code catalogue entry, ``Aldous-Lyons soficity conjecture,'' proposer metadata David Aldous. Dataset curation and metadata: Ulam AI and UnsolvedMath Contributors, CC BY 4.0.
\item David Aldous, \emph{Open problems}, \texttt{unimodular.html}, index \url{https://www.stat.berkeley.edu/~aldous/Research/OP/index.html}; HTML-text extraction with linked primary-PDF text, accessed 2026-07-29.
\item Bowen--Chapman--Lubotzky--Vidick (2024), arXiv:2408.00110, cited in the supplied summary for a nonsofic unimodular random graph; no title is supplied.
\item Bowen--Chapman--Vidick (2024/2025), arXiv:2501.00173, cited there for a soficity-undecidability strengthening; its precise decision problem is not stated.
\item Provenance records 23 active index pages; current primary-literature status audit and redistribution terms NEEDS\_REVIEW, separately from dataset attribution. Neither cited preprint is independently checked here.
\end{itemize}
\end{document}
